% GENERATED FILE. Edit canonical lessons, then run npm run generate:books.
% canonical-source-hash: fnv1a64:7a1ef1575bea0bbc
% canonical-lessons: RU-C285-sobytiye, RU-C285-proisshestviye, RU-C285-khronika, RU-C285-svodka, RU-C285-itog

\chapter{Events and incidents}
\label{ch:ru-events-and-incidents}
\hlchaptermodality{285}

\begin{chapteropening}
\textbf{By the end of this chapter:} \emph{I can talk about what happened: an event, an incident, a news round-up, a summary and the outcome.}

Complete the last lesson of chapter 285: I can talk about what happened: an event, an incident, a news round-up, a summary and the outcome.
\end{chapteropening}

\section[\texorpdfstring{\ru{событие} (sobýtiye)}{событие (sobýtiye)}]{\ru{событие} (sobýtiye) — an event}

\label{lesson:RU-C285-sobytiye}



\begin{quote}
Before the new one: say the Russian for to comment on, then the Russian for to quote.
\end{quote}

\subsection*{You'll want to know: \ru{событие}}
\textbf{\ru{событие}} (\emph{sobýtiye}) — \textquotedblleft{}an event\textquotedblright{}.

\begin{quote}
\ru{Это} \ru{было} \ru{важное} \ru{событие} \ru{для} \ru{всей} \ru{семьи}.
\end{quote}

\begin{quote}
\emph{Éto býlo vázhnoye sobýtiye dlya vsey semí.}
\end{quote}

\begin{quote}
It was an important event for the whole family.
\end{quote}

Neuter. \emph{\ru{Место} \ru{событий}} — the scene.

\begin{cousinweb}
\emph{\ru{событие}} is built on \emph{\ru{со}-} \textquotedblleft{}together\textquotedblright{} and \emph{\ru{бытие}} \textquotedblleft{}being\textquotedblright{}, from \emph{\ru{быть}} \textquotedblleft{}to be\textquotedblright{}, inherited from Proto-Slavic \emph{byti}, from PIE \emph{b\textsuperscript{h}uH-} (EtymDB): what came to be.
\end{cousinweb}

\begin{grammarlens}[title={What you've built}]
One more word for telling, explaining and describing at B1.
\end{grammarlens}

\subsection*{Your turn}
\begin{itemize}
\raggedright
  \item \emph{Say it:} \emph{sobýtiye}
  \item \emph{Say it:} \emph{sobýtiye}, once more
  \item \emph{Say it:} a sentence of your own with \emph{sobýtiye}
  \item \emph{Recall:} say the Russian for to comment on, then the Russian for to quote, then say \emph{sobýtiye} again
\end{itemize}

\begin{tcolorbox}[breakable,colback=teal!4,colframe=teal!35!black,title={Before you move on}]
What does \ru{событие} mean? (\textquotedblleft{}An event\textquotedblright{}.) Say it once more.
\end{tcolorbox}
\section[\texorpdfstring{\ru{происшествие} (proisshéstviye)}{происшествие (proisshéstviye)}]{\ru{происшествие} (proisshéstviye) — an incident}

\label{lesson:RU-C285-proisshestviye}



\begin{quote}
Before the new one: say the Russian for to quote, then the Russian for an event.
\end{quote}

\subsection*{You'll want to know: \ru{происшествие}}
\textbf{\ru{происшествие}} (\emph{proisshéstviye}) — \textquotedblleft{}an incident\textquotedblright{}.

\begin{quote}
\ru{Вчера} \ru{на} \ru{улице} \ru{было} \ru{странное} \ru{происшествие}.
\end{quote}

\begin{quote}
\emph{Vcherá na úlitse býlo stránnoye proisshéstviye.}
\end{quote}

\begin{quote}
Yesterday there was a strange incident in the street.
\end{quote}

Neuter. Usually something unpleasant or unexpected.

\begin{cousinweb}
\emph{\ru{происшествие}} is built on \emph{\ru{произойти}} \textquotedblleft{}to happen\textquotedblright{} with the stem \emph{\ru{шеств}-} \textquotedblleft{}going\textquotedblright{}, seen in Church Slavonic \emph{šĭstvie} \textquotedblleft{}a going, a procession\textquotedblright{} (EtymDB).
\end{cousinweb}

\begin{grammarlens}[title={What you've built}]
One more word for telling, explaining and describing at B1.
\end{grammarlens}

\subsection*{Your turn}
\begin{itemize}
\raggedright
  \item \emph{Say it:} \emph{proisshéstviye}
  \item \emph{Say it:} \emph{proisshéstviye}, once more
  \item \emph{Say it:} a sentence of your own with \emph{proisshéstviye}
  \item \emph{Recall:} say the Russian for to quote, then the Russian for an event, then say \emph{proisshéstviye} again
\end{itemize}

\begin{tcolorbox}[breakable,colback=teal!4,colframe=teal!35!black,title={Before you move on}]
What does \ru{происшествие} mean? (\textquotedblleft{}An incident\textquotedblright{}.) Say it once more.
\end{tcolorbox}
\section[\texorpdfstring{\ru{хроника} (khrónika)}{хроника (khrónika)}]{\ru{хроника} (khrónika) — a news round-up; a chronicle}

\label{lesson:RU-C285-khronika}



\begin{quote}
Before the new one: say the Russian for an event, then the Russian for an incident.
\end{quote}

\subsection*{You'll want to know: \ru{хроника}}
\textbf{\ru{хроника}} (\emph{khrónika}) — \textquotedblleft{}a news round-up; a chronicle\textquotedblright{}.

\begin{quote}
\ru{Каждый} \ru{вечер} \ru{он} \ru{слушает} \ru{хронику} \ru{новостей}.
\end{quote}

\begin{quote}
\emph{Kázhdyy vécher on slúshayet khróniku novostéy.}
\end{quote}

\begin{quote}
Every evening he listens to the news round-up.
\end{quote}

Feminine. \emph{\ru{Хроника}} — a record of events in order.

\begin{cousinweb}
\emph{\ru{хроника}} came through German \emph{Chronik} from Latin \emph{chronica} \textquotedblleft{}a chronicle\textquotedblright{}, from Greek \emph{khronikós} \textquotedblleft{}of time\textquotedblright{}, built on \emph{khrónos} \textquotedblleft{}time\textquotedblright{} (EtymDB).
\end{cousinweb}

\begin{grammarlens}[title={What you've built}]
One more word for telling, explaining and describing at B1.
\end{grammarlens}

\subsection*{Your turn}
\begin{itemize}
\raggedright
  \item \emph{Say it:} \emph{khrónika}
  \item \emph{Say it:} \emph{khrónika}, once more
  \item \emph{Say it:} a sentence of your own with \emph{khrónika}
  \item \emph{Recall:} say the Russian for an event, then the Russian for an incident, then say \emph{khrónika} again
\end{itemize}

\begin{tcolorbox}[breakable,colback=teal!4,colframe=teal!35!black,title={Before you move on}]
What does \ru{хроника} mean? (\textquotedblleft{}A news round-up; a chronicle\textquotedblright{}.) Say it once more.
\end{tcolorbox}
\section[\texorpdfstring{\ru{сводка} (svódka)}{сводка (svódka)}]{\ru{сводка} (svódka) — a summary, a bulletin}

\label{lesson:RU-C285-svodka}



\begin{quote}
Before the new one: say the Russian for an incident, then the Russian for a news round-up.
\end{quote}

\subsection*{You'll want to know: \ru{сводка}}
\textbf{\ru{сводка}} (\emph{svódka}) — \textquotedblleft{}a summary, a bulletin\textquotedblright{}.

\begin{quote}
\ru{По} \ru{радио} \ru{передавали} \ru{сводку} \ru{погоды}.
\end{quote}

\begin{quote}
\emph{Po rádio peredaváli svódku pogódy.}
\end{quote}

\begin{quote}
The radio gave the weather report.
\end{quote}

\emph{\ru{Сводка} \ru{погоды}} — the weather report. Genitive plural \emph{\ru{сводок}}.

\begin{cousinweb}
\emph{\ru{сводка}} is built on \emph{\ru{сводить}} \textquotedblleft{}to bring together\textquotedblright{}, from \emph{\ru{водить}} \textquotedblleft{}to lead\textquotedblright{}, inherited from Proto-Slavic \emph{voditi} (EtymDB): facts brought together.
\end{cousinweb}

\begin{grammarlens}[title={What you've built}]
One more word for telling, explaining and describing at B1.
\end{grammarlens}

\subsection*{Your turn}
\begin{itemize}
\raggedright
  \item \emph{Say it:} \emph{svódka}
  \item \emph{Say it:} \emph{svódka}, once more
  \item \emph{Say it:} a sentence of your own with \emph{svódka}
  \item \emph{Recall:} say the Russian for an incident, then the Russian for a news round-up, then say \emph{svódka} again
\end{itemize}

\begin{tcolorbox}[breakable,colback=teal!4,colframe=teal!35!black,title={Before you move on}]
What does \ru{сводка} mean? (\textquotedblleft{}A summary, a bulletin\textquotedblright{}.) Say it once more.
\end{tcolorbox}
\section[\texorpdfstring{\ru{итог} (itóg)}{итог (itóg)}]{\ru{итог} (itóg) — the result, the total}

\label{lesson:RU-C285-itog}



\begin{quote}
Before the new one: say the Russian for a news round-up, then the Russian for a summary.
\end{quote}

\subsection*{You'll want to know: \ru{итог}}
\textbf{\ru{итог}} (\emph{itóg}) — \textquotedblleft{}the result, the total\textquotedblright{}.

\begin{quote}
\ru{В} \ru{итоге} \ru{мы} \ru{решили} \ru{остаться} \ru{дома}.
\end{quote}

\begin{quote}
\emph{V itóge my reshíli ostátsya dóma.}
\end{quote}

\begin{quote}
In the end we decided to stay at home.
\end{quote}

\emph{\ru{В} \ru{итоге}} — in the end; \emph{\ru{подвести} \ru{итоги}} — to sum up.

\begin{cousinweb}
\emph{\ru{итог}} is usually traced to \emph{\ru{и}-}, a form of \emph{\ru{из}} \textquotedblleft{}out\textquotedblright{}, and the root of \emph{\ru{течь}} \textquotedblleft{}to flow\textquotedblright{}: what flows out at the end. EtymDB lists no parent.
\end{cousinweb}

\begin{grammarlens}[title={What you've built}]
One more word for telling, explaining and describing at B1.
\end{grammarlens}

\subsection*{Your turn}
\begin{itemize}
\raggedright
  \item \emph{Say it:} \emph{itóg}
  \item \emph{Say it:} \emph{itóg}, once more
  \item \emph{Say it:} a sentence of your own with \emph{itóg}
  \item \emph{Recall:} say the Russian for a news round-up, then the Russian for a summary, then say \emph{itóg} again
\end{itemize}

\begin{tcolorbox}[breakable,colback=teal!4,colframe=teal!35!black,title={Before you move on}]
What does \ru{итог} mean? (\textquotedblleft{}The result, the total\textquotedblright{}.) Say it once more.
\end{tcolorbox}
